\documentclass{package/fancy-book}
\setlength{\parindent}{2em}
\usepackage[UTF8]{ctex}

%%%%%%%%%% Default Package %%%%%%%%%%%%%
\usepackage{package/color-env}
\usepackage{package/quiver}
\usepackage{background}
\usepackage[object=vectorian]{pgfornament} %% used in title.tex
\usepackage{calligra} %%% (optional) to make the Title text beautiful 
\usepackage{lipsum}  %% for dummy text 
\usepackage{amssymb,amsmath,amsfonts}  %%% for maths
\usepackage{datetime}
%%%%%% Optional Packages %%%%%%%
\usepackage{lettrine} %% for nice looking 
\usepackage{GoudyIn} %% first Letter of the paragraph
\renewcommand{\LettrineFontHook}{\color{black}\GoudyInfamily{}}
\LettrineTextFont{\itshape}
\setcounter{DefaultLines}{3}%
%%%%%%%%%%%%%%%%%%%%%%%%%%%%%%%%%%%%%
\usepackage{datetime2}
\usepackage{fourier-orns}
\newcommand{\ornamento}{\vspace{2em}\noindent \textcolor{darkgray}{\hrulefill~ \raisebox{-2.5pt}[10pt][10pt]{\leafright \decofourleft 
\decothreeleft  \aldineright \decotwo \floweroneleft \decoone   \floweroneright 
\decotwo \aldineleft\decothreeright \decofourright \leafleft} ~  \hrulefill \\ \vspace{2em}}}

%定义命令
\newcommand{\N}{\mathbb{N}}
\newcommand{\Z}{\mathbb{Z}}
\newcommand{\Q}{\mathbb{Q}}
\newcommand{\R}{\mathbb{R}}
\newcommand{\C}{\mathbb{C}}
\newcommand{\HH}{\mathbb{H}}
\newcommand{\RP}{\mathbb{RP}}

%%%% Bibliography %%%%%%%%%
% Required packages are included in notes class
% Can be tweaked in the notes.cls file itself
\addbibresource{resource/references.bib}

\begin{document}
\begin{titlepage}

%%%%%%%%%%%%%%%%%%%%%%%%%%%%%%%%%%%% Inspired From  %%%%%%%%%%%%%%%%%%%%%%%%%%%%%%%%%%%%%%%%%%%%%%%%% 
%%%%  https://www.reddit.com/r/LaTeX/comments/j9d739/hello_world_in_latex_is_a_lot_cooler/   %%%%%
%%%%%%%%%%%%%%%% If this doesn't look nice then you may remove it %%%%%%%%%%%%%%%%%%%%%%%%%%

\backgroundsetup{
scale=1,
opacity=1,
angle=0,
color=black,
contents={
\begin{tikzpicture}[color=black, every node/.style={inner sep= 15pt}]
\node (NW) [anchor=north west] at (current page.north west){\pgfornament[width=2.5cm] {131}};
\node (NE) [anchor=north east] at (current page.north east){\pgfornament[width=2.5cm, symmetry=v]{131}};
\node (SW) [anchor=south west] at (current page.south west){\pgfornament[width=2.5cm, symmetry=h]{131}};
\node (SE) [anchor=south east] at (current page.south east){\pgfornament[width=2.5cm, symmetry=c]{131}};
\foreach \i in {-4,0,4}
\node[anchor=north,xshift=\i cm] at (current page.north){\pgfornament[scale=0.25,symmetry=v]{71}};
\foreach \i in {-4,0,4}
\node[xshift=\i cm, yshift=32.25 pt] at (current page.south){\pgfornament[scale=0.25,symmetry=v]{71}};
\foreach \i in {-8,-4,0,4,8}
\node[yshift=\i cm, xshift=32.25pt, rotate=90] at (current page.west){\pgfornament[scale=0.25,symmetry=v]{71}};
\foreach \i in {-8,-4,0,4,8}
\node[yshift=\i cm, xshift=-32.25pt, rotate=90] at (current page.east){\pgfornament[scale=0.25,symmetry=v]{71}};
\foreach \i in {-11,-9,...,7,9}
\node[anchor=west, yshift=\i cm, xshift=52.25pt, rotate=90] at (current page.west){\pgfornament[scale=0.1]{80}};
\foreach \i in {-11,-9,...,7,9}
\node[anchor=east, yshift=\i cm, xshift=-52.25pt, rotate=-90] at (current page.east){\pgfornament[scale=0.1]{80}};
\end{tikzpicture}
}}

%%%%%%%%%%%%%%%%%%%%%%%%%%%%%%%%%%%%%%%%%%%%%%%%%%%%%%%%%%%%%%%%%%%%%%%%%%%%%%%%%%%%%%%%%%%%%%%%%%%%%%%%%%%%%%%%%%

\centering % Centre everything on the title page
		
\scshape % Use small caps for all text on the title page

\vspace*{\baselineskip} % White space at the top of the page

%------------------------------------------------
%	Title
%------------------------------------------------

\rule{\textwidth}{1.6pt}\vspace*{-\baselineskip}\vspace*{2pt} % Thick horizontal rule
\rule{\textwidth}{0.4pt} % Thin horizontal rule

\vspace{0.75\baselineskip} % Whitespace above the title

{\huge \calligra{ \textbf{} }\\} % Title

\vspace{0.75\baselineskip} % Whitespace below the title

\rule{\textwidth}{0.4pt}\vspace*{-\baselineskip}\vspace{3.2pt} % Thin horizontal rule
\rule{\textwidth}{1.6pt} % Thick horizontal rule

\vspace{2\baselineskip} % Whitespace after the title block

%------------------------------------------------
%	Subtitle
%------------------------------------------------

\LARGE{Differential Topology} 

\vspace*{3\baselineskip} % Whitespace under the subtitle



\vspace{0.5\baselineskip} 

{\scshape   \LARGE Edited by\\  颜成子游/南郭子綦} % Editor list
\vspace{0.2\baselineskip} 


\vfill 
{最后一次编译时间：\DTMnow}
%------------------------------------------------
% Author
%------------------------------------------------

\begin{figure}[!h]
    \centering
    \includegraphics[width = 3cm, height= 3cm]{resource/icon.png}%% include the university icon here
\end{figure}
\vspace{0.3\baselineskip} 

\end{titlepage}
\backgroundsetup{contents={}} %% to remove background and watermark from other pages
\tableofcontents

\quad


\chapter{De Rham理论}
\section{}
\begin{proposition}{}
   
\end{proposition}
\chapter{\u{C}ech-de Rham复形}

\chapter{代数拓扑与谱序列}

\chapter{示性类}
\section{}
\begin{theorem}[Borsuk-Ulam定理]{Borsuk-Ulam定理1}
    设$f:S^n \to \mathbb{R}^n$是连续函数。则一定存在点$p \in S^n$使得$f(p)=f(-p)$。
\end{theorem}
\end{document}